\section{The determinant and integration by parts}\label{sec:calculus}

The wave current in Proposition~\ref{prop:current} is positive and has
divergence-free real columns. We shall pair it with a Hessian. The
essential pairing is the ordinary real contraction
$A:H(v)=A^{\RR}:D^2v$: each spacetime index still has its three real
coordinates. The choice of Hessian has two further consequences. The squared modulus of a
holomorphic function contributes at most one quaternionic direction.
In $n$ spatial variables, if its nonzero eigenvalue is $\beta$, adding
$\epsilon I$ gives determinant $(\beta+\epsilon)\epsilon^{n-1}$.
This is where the denominator in $p_{\mathrm c}=2n/(n-1)$ enters.
The cofactor of the Hessian also has
divergence-free columns; this allows us to integrate the determinant
equation by parts, including at the boundary.

\subsection{The Hessian in real and complex coordinates}

Use the triples $z_a=x_a+P_a\mathbf i\in\mathcal D$ of
Section~\ref{sec:current}, with $P_a\in\CC_{\mathbf k}$. We allow a
fourth real coordinate as well: Appendix~\ref{sec:dirichlet} uses it to
rotate the normal at the boundary. Thus for each of $m$ coordinate
indices let $z_a=\sum_{\alpha=0}^{k-1} e_\alpha z_{a\alpha}$, where
$k=3$ or $4$, $(e_0,e_1,e_2,e_3)=(1,\mathbf i,\mathbf j,\mathbf k)$,
and the $z_{a\alpha}$ are real.
For a real $C^2$ function $u$ put $\partial_{a\alpha}=\partial/\partial z_{a\alpha}$
and define
\begin{equation}\label{eq:quaternionic-hessian}
  G_a(u)=\sum_\alpha e_\alpha\,\partial_{a\alpha}u,\qquad
  H_{ab}(u)=\sum_{\alpha,\beta} e_\alpha\bar e_\beta\,
                        \partial_{a\alpha}\partial_{b\beta}u .
\end{equation}
Thus $G(u)$ is the ordinary gradient written as a quaternionic column.
The symbols $\nabla$, $\Delta$ and $\partial_{\mathbf n}$ denote the
Euclidean gradient, Laplacian and outward normal derivative.

\begin{lemma}[Hessian calculus]\label{lem:hessian-calculus}
In either number of slots, $H(u)$ is Hermitian and $\tr H(u)=\Delta u$.
The linear map from the real Hessian to $H(u)$ sends
$\xi\xi^{\mathsf t}$ to $gg^*$, where
$g_a=\sum_\alpha e_\alpha\xi_{a\alpha}$. It sends nonnegative
real forms to nonnegative quaternionic forms. The product rule is
\begin{equation}\label{eq:hessian-product}
 H(uv)=uH(v)+vH(u)+G(u)G(v)^*+G(v)G(u)^* .
\end{equation}
In four slots $H$ is covariant under quaternionic unitary changes of
coordinates, and in three slots under real orthogonal changes of the
coordinate index applied in all slots.

If $u$ depends only on $P\in\CC_{\mathbf k}^m$, then
\begin{equation}\label{eq:complex-hessian}
 G_a(u)=2(\partial_{\bar P_a}u)\mathbf i,\qquad
 H_{ab}(u)=4\partial_{\bar P_a}\partial_{P_b}u .
\end{equation}
In particular, a holomorphic map $F=(F_1,\ldots,F_l)$ with complex
Jacobian $J_{\nu a}=\partial_{P_a}F_\nu$ satisfies
\begin{equation}\label{eq:holomorphic-squares}
 H\Big(\sum_\nu|F_\nu|^2\Big)=4J^*J .
\end{equation}
For every smooth real function $\rho$ of $F$,
$G(\rho\circ F)$ belongs to the quaternionic image of $J^*$.
\end{lemma}

\begin{proof}
Interchanging the real derivatives gives Hermitian symmetry. In a
diagonal entry the terms with $\alpha\ne\beta$ cancel in pairs; the
remaining terms have coefficient one. Formula~\eqref{eq:quaternionic-hessian}
applied to $\xi\xi^{\mathsf t}$ gives $gg^*$, and sums of these
rank-one forms give every nonnegative real form. The same formula and
the ordinary product rule give \eqref{eq:hessian-product}.
In four slots the real transpose of a quaternionic unitary $U$ acts
on the encoded covector as $U^*$. Thus the Hessian rule for a real
rank-one form gives $H(u\circ U)=U^*(H(u)\circ U)U$; linearity
extends this to every real Hessian. Real orthogonal matrices on the
coordinate index also preserve the three-slot subspace, giving the
stated covariance there.

For $P_a=p_a+\mathbf k q_a$ we have
$2\partial_{\bar P_a}=\partial_{p_a}+\mathbf k\partial_{q_a}$.
Since $\mathbf k\mathbf i=\mathbf j$, this gives the gradient in
\eqref{eq:complex-hessian}. The four terms in $H_{ab}$ have
coefficients $1,-\mathbf k,\mathbf k,1$ in the order
$(p_ap_b,p_aq_b,q_ap_b,q_aq_b)$, which gives its Hessian formula.
The complex chain rule now gives \eqref{eq:holomorphic-squares} and
$G(\rho\circ F)=2J^*(\partial_{\bar F_\nu}\rho)_\nu\,\mathbf i$,
with the derivatives of $\rho$ evaluated at $F$.
\end{proof}

\subsection{The Moore cofactor}

For Hermitian quaternionic matrices use the real inner product
$A:B=\Rea\tr(AB)$. The real part of the trace is cyclic; in particular
$A:(vv^*)=v^*Av$. Write $\det A$ for the Moore determinant, the product
of the $m$ real eigenvalues counted with quaternionic multiplicity.
These conventions agree with \cite[Section~1.1, especially Theorems
1.1.8--1.1.9]{Alesker2003LA}. The next identities apply also to the
full Hessian in the energy argument, whose spatial block is positive
definite but whose time entry need not have a sign.

\begin{lemma}[Matrix identities]\label{lem:quaternionic-matrices}
The Moore determinant is a homogeneous polynomial of degree $m$.
Its cofactor $C(A)$ is the Hermitian matrix defined by
\begin{equation}\label{eq:cofactor-definition}
 C(A):B=\left.\frac{d}{dt}\det(A+tB)\right|_{t=0}.
\end{equation}
It is polynomial of degree $m-1$, and
\begin{equation}\label{eq:cofactor-inverse}
 C(A)=(\det A) A^{-1}\quad (A>0).
\end{equation}
If $A\ge0$, then $C(A)\ge0$; if also $A$ is singular, $C(A)$ kills
the image of $A$. On positive definite matrices $\log\det$ is concave
and increasing in the matrix order.

For a unit vector $v\in\HH^m$ let $A|_{v^\perp}$ denote the
compression of the form to its quaternionic orthogonal complement.
For every Hermitian $A$ and real $t$,
\begin{equation}\label{eq:cofactor-restriction}
  v^*C(A)v=\det(A|_{v^\perp}),\qquad
  \det(A+t vv^*)=\det A+t\,v^*C(A)v.
\end{equation}
The determinant on the zero-dimensional space is one. Finally, for a
Hermitian $M=\bigl(\begin{smallmatrix}\ell&c^*\\ c&D\end{smallmatrix}\bigr)$
with $D>0$, put $\sigma=\ell-c^*D^{-1}c$ and let $v_0$ be the first
coordinate vector. Then
\begin{equation}\label{eq:schur-complement}
     \det M=(\det D)\sigma,\qquad M-\sigma v_0v_0^*\ge0 .
\end{equation}
\end{lemma}

\begin{proof}
A Hermitian matrix has a quaternionic orthonormal eigenbasis:
maximize $v^*Av$ on the unit sphere and repeat on the orthogonal
complement. By Newton's identities its determinant is a polynomial
in the power sums $\Rea\tr(A^j)$, $1\le j\le m$. This gives
\eqref{eq:cofactor-definition} with the stated degrees.
In the real representation on $\RR^{4m}$ each eigenvalue occurs
four times. Thus for $A>0$, $\log\det A$ is one fourth of the real
log determinant. Differentiation proves \eqref{eq:cofactor-inverse},
monotonicity and concavity. On a diagonal matrix each cofactor
eigenvalue is the product of the other $m-1$ eigenvalues, first for
positive matrices by \eqref{eq:cofactor-inverse}, then for all
diagonal matrices by polynomiality. Unitary invariance proves the
claims for singular $A$.

Completing the square shows that the form of $M-\sigma v_0v_0^*$ at
$(b,w)$ is $(w+D^{-1}cb)^*D(w+D^{-1}cb)$. For $\sigma>0$ the real
Schur determinant formula gives $(\det M)^4=(\det D)^4\sigma^4$;
take the positive fourth root. Both sides of $\det M=(\det D)\sigma$
are polynomials in $\ell$, so the identity holds for every $\sigma$.
In a basis with $v$ first, adding $tvv^*$ changes only $\ell$.
For positive $A$ the change in its determinant is therefore exactly
$t\det(A|_{v^\perp})$. Differentiation in $t$ gives the first formula
in \eqref{eq:cofactor-restriction}; polynomiality extends both to
every Hermitian $A$. For $m=1$ they hold directly.
\end{proof}

\subsection{Column divergence and the boundary term}

In $k$ slots the real tensor \eqref{eq:current-real-tensor} of a
Hermitian field $A(z)$ has coefficients
\begin{equation}\label{eq:real-coefficients}
     A^{\RR}_{a\alpha,b\beta}=\Rea(\bar e_\alpha A_{ab}e_\beta).
\end{equation}
This symmetric real matrix represents the Hermitian form on the
encoded real vectors: $(\nabla v)^{\mathsf t} A^{\RR}\nabla u=
\Rea(G(v)^*A G(u))$. Likewise $A:H(u)=A^{\RR}:D^2u$, where the
pairing on the right is the ordinary real matrix inner product.
Consequently, whenever the real columns have zero divergence,
\begin{equation}\label{eq:real-divergence}
    A:H(u)=\Div(A^{\RR}\nabla u).
\end{equation}
One sufficient condition is the quaternionic row identity
\begin{equation}\label{eq:strong-divergence}
     \sum_{a,\alpha}\bar e_\alpha\,\partial_{a\alpha}A_{ab}=0
           \quad(1\le b\le m):
\end{equation}
multiply on the right by $e_\beta$ and take real parts. Also,
if $A G(w)=0$ for a real scalar function $w$, then $wA$ preserves
either form of column conservation already satisfied by $A$.
Indeed the extra row in the product rule is $G(w)^*A=0$.

Alesker showed that the real tensor representing the Moore cofactor
of a positive definite quaternionic Hessian has divergence-free columns.
This puts the linearized Monge--Ampère operator in divergence form
\cite[Lemma~4.6]{Alesker2013Flat}. The following polynomial proof
also covers the three slots used by the wave current and allows an
indefinite or singular Hessian.

\begin{lemma}[Cofactor divergence]\label{lem:cofactor-divergence}
For $k=3$ or $4$ and a smooth real $u$, the Hessian cofactor satisfies
\begin{equation}\label{eq:cofactor-divergence}
    \sum_{a,\alpha}\bar e_\alpha\,\partial_{a\alpha} C(H(u))_{ab}=0
    \quad(1\le b\le m).
\end{equation}
\end{lemma}

\begin{proof}
For $m=1$ the cofactor is one. Otherwise put $j_0=m-1$. Polarize
the cofactor: for Hermitian inputs $M_1,\ldots,M_{j_0}$, denote
the coefficient of $t_1\cdots t_{j_0}$ in $C(\sum_j t_jM_j)$ by
$Q(M_1,\ldots,M_{j_0})$. This is multilinear, and on the diagonal
equals $j_0!\,C(M_1)$.

Take $M_j=g_jg_j^*$. For $t_j>0$, $\sum_jt_jg_jg_j^*$ is singular:
its image is the quaternionic span of the $j_0$ columns. Its cofactor
therefore kills each $g_j$, by Lemma~\ref{lem:quaternionic-matrices}.
This polynomial identity in $t$ holds coefficient by coefficient.
Thus $Q=Q(g_1g_1^*,\ldots,g_{j_0}g_{j_0}^*)$ kills every $g_j$,
and $g_j^*Q=0$ as well since $Q$ is Hermitian.

Test the polarized identity on $u_j(z)=\exp(\xi_j\cdot z)$ with
arbitrary real covectors $\xi_j$, using the Euclidean dot product.
Set $(g_j)_a=\sum_\alpha e_\alpha(\xi_j)_{a\alpha}$. The Hessian
symbol in Lemma~\ref{lem:hessian-calculus} gives
$Q(H(u_1),\ldots,H(u_{j_0}))=\exp((\sum_j\xi_j)\cdot z)\,Q$.
After cancellation of this exponential, its row divergence is
$(\sum_jg_j)^*Q=0$. To pass to arbitrary smooth inputs, observe
that the row divergence is a multilinear differential expression with
constant coefficients: its value on these exponential functions,
divided by their product, is the polynomial whose coefficients
multiply the corresponding derivatives of the inputs. That polynomial
vanishes for all covectors, so all coefficients vanish. Taking every
input equal to $u$ proves \eqref{eq:cofactor-divergence}.
\end{proof}

Here is the boundary convention used for both the current and the
cofactor. On a smooth domain $\Omega=\{r<0\}$ assume $\nabla r\ne0$
on $\partial\Omega$. Write $\mathbf n=\nabla r/|\nabla r|$ for the
real outward unit normal and $\nu=G(r)/|\nabla r|$ for the same
vector in quaternionic coordinates. By \eqref{eq:real-divergence},
Green's identity for a field with divergence-free real columns is
\begin{equation}\label{eq:calculus-green}
 \int_\Omega\!\big(v\,A:H(u)-u\,A:H(v)\big)
  =\int_{\partial\Omega}\!
       \Rea\!\left(v\,\nu^*A G(u)-u\,\nu^*A G(v)\right)\dd S .
\end{equation}
Volume and surface measure are Euclidean. The same formula holds in
periodic coordinates, with test functions of compact support in any
unbounded directions.

If $u$ vanishes on $\partial\Omega$, then $G(u)=\nu\,\partial_{\mathbf n}u$
there. Also $u=\rho r$ for a smooth $\rho$ near the boundary: use $r$
as one local coordinate and integrate the derivative of $u$ in that
coordinate. The cross
terms of \eqref{eq:hessian-product} then vanish on compression, giving
\begin{equation}\label{eq:boundary-restriction}
 H(u)|_{\nu^\perp}=\rho\,H(r)|_{\nu^\perp}\qquad\text{on }\partial\Omega .
\end{equation}
In four slots this calculation gives the invariance of the quaternionic
Levi form under a change of defining function \cite{Alesker2003MA}.
Functions independent of the fourth slot give the same formula in three slots.
In particular the normal flux coefficient $\nu^*C(H(u))\nu$ depends
only on this compressed form, by \eqref{eq:cofactor-restriction}.
This is why the determinant boundary term can be evaluated without
knowing the second derivative of $u$ in the real normal direction.
